% ------------------------------------------------------------
% CHAPTER 32
% ------------------------------------------------------------

\chapter{Continuous-Time ΔΣ and Field-Theoretic Limits}

The discrete-time delta–sigma modulator embodies a carefully engineered balance between quantization, feedback, and stability. At each sampling step, it compresses the continuum of input states into a discrete symbol, injects entropy into the internal reservoir, and employs controlled feedback flows to route this entropy into spectrally remote regions. When the sampling frequency greatly exceeds the signal bandwidth, the modulator behaves approximately as a continuous dynamical system punctuated by discrete collapse events. In the limit of infinite oversampling, these punctuations become infinitesimally spaced, giving rise to a continuous-time model whose behavior resembles a field theory defined on a one-dimensional time manifold. This chapter develops the mathematical framework for this limiting object, connecting the classical continuous-time delta–sigma architecture with the entropy-based and geometric interpretations introduced earlier.

To begin, consider a continuous-time model with an internal state \(u(t)\), an input signal \(\Phi(t)\), and a quantizer that produces a piecewise-constant output \(y(t)\) that tracks the sign or region of the internal state. The evolution equation for the internal variable takes the form
\[
\frac{d u(t)}{dt} = \mathcal{F}(\Phi(t), y(t)),
\]
where \(\mathcal{F}\) encapsulates the loop filter. In first-order modulators,
\[
\frac{d u(t)}{dt} = \Phi(t) - y(t).
\]
Higher-order modulators introduce additional states \(u_1(t), u_2(t)\), and so forth, each governed by successive integrations and weighted combinations of the input and quantized output. The quantizer operates at discrete instants and determines the value of \(y(t)\) for the intervening interval. When the sampling rate tends to infinity, the quantized output becomes a rapidly switching signal approximating a distributional object. The internal dynamics must then be interpreted as differential inclusions rather than classical differential equations.

A central challenge arises from the discontinuity of the quantizer. Even in the continuous-time limit, the symbol collapse map remains discontinuous, which means that the classical theory of smooth differential equations is insufficient for describing the modulator’s behavior. Instead, one must adopt the language of Filippov differential inclusions. In this framework, the delta–sigma loop becomes a system whose right-hand side is multivalued along the quantization manifold. The differential inclusion for the first-order modulator is
\[
\frac{d u(t)}{dt} \in \Phi(t) - \mathrm{Conv}\{y^+, y^-\},
\]
where \(y^+\) and \(y^-\) are the quantized values immediately above and below the discontinuity, and the convex hull yields the set of admissible velocities. This formalism makes it possible to interpret the modulator's fast-switching behavior as motion along a sliding mode surface of the inclusion.

When oversampling tends to infinity, the bitstream becomes so rapidly oscillatory that its average value converges to a continuous function. In the first-order modulator, the average of the output converges to the average of the input. In higher-order modulators, the same convergence holds after removing the filtering induced by the internal dynamics. In all cases, the continuous-time model exhibits a form of averaging that allows the discontinuous quantized output to approximate a smooth trajectory. Formally, if \(y_\varepsilon(t)\) denotes the output of a modulator operating at oversampling factor \(\varepsilon^{-1}\), then under appropriate stability assumptions,
\[
\lim_{\varepsilon \to 0} y_\varepsilon(t) = \bar{y}(t),
\]
where \(\bar{y}(t)\) satisfies a differential equation determined by the loop filter and the desired signal-reconstruction properties. This convergence justifies representing the modulator as a continuous-time system equipped with an effective nonlinearity replacing the quantizer.

To incorporate thermodynamics, consider the entropy associated with the internal state distribution. In continuous time, the reservoir state \(u(t)\) develops a time-dependent probability density \(\rho(u,t)\) influenced by the deterministic flow and by the stochastic microstructure generated by rapid quantization events. When the oversampling tends to infinity, the randomness associated with discrete collapse becomes diffusive. One obtains a Fokker–Planck-like evolution equation of the form
\[
\frac{\partial \rho}{\partial t} = - \nabla_u \cdot \big( \rho\, v(u,t) \big) + D \nabla_u^2 \rho,
\]
where \(v(u,t)\) is the drift induced by the loop filter and the effective nonlinearity, and \(D\) is a diffusion coefficient proportional to the rate at which quantization events perturb the internal state. The delta–sigma modulator therefore approaches a field theory for the internal state distribution: the internal dynamics are described by a drift–diffusion equation that captures both the deterministic action of the loop filter and the stochastic effect of quantization in the continuous limit.

Entropy production in this regime is governed by
\[
\frac{d S(t)}{dt} = \int \rho(u,t) \nabla_u \cdot v(u,t) \, du + D \int \frac{|\nabla_u \rho|^2}{\rho} \, du,
\]
revealing that the divergence of the drift contributes to entropy change, while the diffusion term always increases entropy. A well-designed modulator enforces negative divergence in the drift field over the signal band, ensuring that entropy does not accumulate in the region of interest. Instead, positive divergence appears in the high-frequency region of the output where entropy can be harmlessly expelled. The thermodynamic principle governing the continuous-time delta–sigma loop is therefore that entropy must be routed away from the degrees of freedom relevant for reconstruction.

The connection to RSVP field theory becomes explicit when we interpret \(u(t)\) not as a scalar but as a field value. Suppose the signal is a spatially extended object \(\Phi(x,t)\) defined over some domain \(x \in \Omega\). A distributed delta–sigma modulator can be defined that samples the field at each point and generates local collapse events. In the continuous-time limit, both the collapse and the feedback become field operations. The internal field \(U(x,t)\) then satisfies a partial differential equation of the form
\[
\frac{\partial U}{\partial t} = \mathcal{D}[U] + \Phi(x,t) - Y(x,t),
\]
where \(\mathcal{D}\) is a spatial differential operator representing diffusion or wave propagation, and \(Y(x,t)\) is the quantized field obtained by collapsing \(U(x,t)\) according to a local rule. The oversampling limit smooths \(Y(x,t)\) into a field obeying an average equation, thereby transforming the discontinuous collapse dynamics into a continuous PDE.

This field-theoretic limit reveals that delta–sigma modulation is a primitive instance of a broader class of measurement operators acting on fields. The modulator becomes an instrument that enforces a relationship between the unobserved field \(\Phi(x,t)\) and the discrete observations \(Y(x,t)\). The feedback loop ensures that the observed field tracks the unobserved one while routing entropy into spectral or spatial regions where it does not interfere with reconstruction. This is precisely the structure of RSVP, where scalar fields represent potential, vector fields represent flows, and entropy fields represent disorder. The delta–sigma loop becomes a reduced model of RSVP dynamics, with the quantization collapse corresponding to entropy injection and the feedback loop corresponding to the smoothing, routing, and dissipation of entropy.

In conclusion, the continuous-time delta–sigma modulator emerges as a field-theoretic limit of the discrete-time entropy engine described earlier. The discontinuous quantizer becomes an effective nonlinearity described by differential inclusions; the internal state becomes a field obeying drift–diffusion dynamics; and the output becomes a field satisfying an averaged differential equation. The modulator therefore provides a bridge between engineered measurement systems and the general field-theoretic structures that govern entropy flow, coherence, and information in physical and cognitive systems.

The next part of the book begins the exploration of how delta–sigma modulation appears in natural and artificial systems, beginning with biological sensing and culminating in large-scale cognitive and AI architectures.

